\documentclass[runningheads]{llncs}

\begin{document}

\title{PaperScan: Sistema de Recuperacion de Informacion para Investigacion Cientifica y Academica}
\author{Equipo PaperScan}
\institute{Universidad}
\maketitle

\begin{abstract}
Este documento describe la arquitectura y el estado actual del sistema PaperScan,
un proyecto de recuperacion de informacion orientado al dominio de investigacion
cientifica y academica. Se presentan los componentes de la Primera Entrega:
adquisicion, indexacion, recuperacion no basica y base vectorial inicial, usando
como base los PDF locales y con respaldo web en arXiv cuando la recuperacion local
es insuficiente.
\keywords{recuperacion de informacion \and arXiv \and indexacion \and modelo de lenguaje}
\end{abstract}

\section{Introduccion}
PaperScan es un sistema de recuperacion de informacion que obtiene documentos
desde un repositorio local de PDF, construye estructuras de indexacion y aplica
un modelo no basico para recuperar documentos relevantes. Cuando los resultados
locales son insuficientes, activa una busqueda de respaldo en arXiv.

\section{Dominio}
El dominio seleccionado es \textit{investigacion cientifica y academica}, con
fuentes de datos basadas en publicaciones de arXiv.

\section{Arquitectura Modular}
La arquitectura del proyecto se organiza en modulos:
\begin{itemize}
\item Adquisicion de datos.
\item Crawler y scraping web.
\item Indexacion.
\item Recuperacion.
\item Base vectorial.
\item RAG (pendiente).
\item Posicionamiento (pendiente).
\item Busqueda web (respaldo arXiv).
\item Modulos opcionales: expansion y retroalimentacion, multimodal,
recomendacion y evaluacion (pendientes).
\end{itemize}

\section{Implementacion Actual}
\subsection{Adquisicion de datos}
Se construye un corpus local en JSONL a partir de los PDF ubicados en
\texttt{local/local\_papers}. Para cada documento se extrae texto y se generan
campos estructurados para indexacion.

Adicionalmente, se implementa un cliente de arXiv con politicas de crawling
(demora entre solicitudes, tamano de lote y reintentos) para respaldo web.

\subsection{Crawler y scraping web}
Se implementa un crawler web basico con profundidad limitada y respeto a
\texttt{robots.txt}, orientado al dominio de investigacion cientifica y academica.
El crawler descarga HTML, extrae texto limpio, y almacena evidencias en
\texttt{datos/brutos/crawl\_web} junto con un archivo \texttt{registros.jsonl}
para trazabilidad de URLs, profundidad y fecha de descarga.

\subsection{Indexacion}
Se realiza preprocesamiento de texto y construccion de indice invertido,
persistiendo estructuras en disco para uso posterior.

\subsection{Recuperacion}
Se implementa un modelo probabilistico de lenguaje (Query Likelihood) con
suavizado Jelinek-Mercer para ranking inicial de resultados locales.

\subsection{Base vectorial}
Se construyen representaciones TF-IDF para almacenamiento vectorial inicial.

\subsection{Salida de resultados}
El sistema persiste resultados y evidencia de ejecucion en \texttt{datos}:
corpus bruto, indice invertido, vectores, estadisticas del corpus local y
resultados de recuperacion.

\subsection{Estadisticas del corpus local}
Para la evidencia de Corte 1, el corpus local generado reporta:
\begin{itemize}
\item Cantidad de documentos: 50.
\item Promedio de palabras por resumen: 267.44.
\item Categorias mas frecuentes: investigacion\_cientifica (50), academico (50).
\end{itemize}

\section{Flujo Operacional}
\begin{enumerate}
\item Cargar y transformar PDF locales a corpus JSONL.
\item Indexar el corpus y construir base vectorial inicial.
\item Ejecutar recuperacion local (modelo de lenguaje y similitud vectorial).
\item Si la recuperacion local no alcanza un umbral minimo de resultados,
activar respaldo web en arXiv.
\item Guardar respuesta final y resultados en disco.
\end{enumerate}

\section{Ejecucion}
El punto de entrada del sistema es \texttt{main.py}. La descarga local de PDF
se realiza desde \texttt{local/scrapping.py}. La consulta se puede pasar por
linea de comandos con \texttt{--consulta}.

\section{Trabajo Futuro}
Las siguientes fases incorporaran integracion RAG, modulo de posicionamiento,
busqueda web automatica e interfaz visual completa.

\section{Bibliografia Base}
\begin{thebibliography}{1}
\bibitem{croft2010}
Croft, W. B., Metzler, D., Strohman, T.:
Search Engines: Information Retrieval in Practice.
Addison-Wesley (2010)
\end{thebibliography}

\end{document}
